% Sections 4--6 (construction, security, implementation) of the revised OJ-CS manuscript.
\section{Construction}\label{sec:construction}
\subsection{Building blocks}
We work in BLS12-381.
BBS+ signs a vector of messages; a holder proves knowledge of a signature while revealing any subset, and Pedersen commitments allow blind issuance~\cite{au2006bbsplus,camenisch2016bbsplus,pedersen1991}; a VB accumulator holds revoked identifiers; a holder proves that a hidden value is not a member~\cite{vitto2022vb}; Bulletproofs++ prove that a committed value lies in a range~\cite{eagen2024bpp}.
All statements of one operation form a single Fiat--Shamir proof, in which witness equalities force the same hidden value across statements~\cite{dockcrypto}. All bases and proving keys are derived by hashing public labels. This matters: an issuer that chose the accumulator's proving key could strip the randomization of non-membership proofs and recover hidden identifiers.

\subsection{Mandate layout}
A mandate is a BBS+ signature on $L=13+n$ messages ($n=16$ classes, $L=\NMsgs$): $s$, $\mathit{id}$, $\mathit{anc}_{1..4}$, $\delta$, the scope bits $b_{1..n}$, $\beta$, $\ell$, $k$, $\tau$, $t$ and a key version. All mandates of a tree share $s$ and $t$. The issuer chooses $\mathit{id}$; $s$ and $t$ are never shown to it.

\subsection{Algorithms}
\paragraph*{IssueRoot} The wallet draws $s,t$ and sends $C=h_0^{r}h_s^{s}h_t^{t}$, $T=g^{t}$, a Schnorr signature by $u_P$ on (label, $T$), and a proof that $C$ commits to $\log_g T$, bound to an issuer nonce. $\mathcal{I}$ authenticates $P$ and checks the proof, the signature, that $T$ is new, the layout of the signed vector, the bounds $\beta,k,\tau<2^{40}$ and $k-\ell\le Q$, and its per-principal root limit. It then draws a fresh non-zero $\mathit{id}$, blind-signs, and logs the binding. The wallet fixes every clear attribute ($\ell_0,k_0,\beta,\tau$, scope) in its request and checks that the returned signature is on exactly those values.

\paragraph*{Present} An invocation $\iota=(\theta,\mathsf{call},c,\mathit{amt},\mathit{cur},e,\mathit{tm},n_\theta)$ consists of the tool's public identity $\theta$, the canonical call, its class, amount and currency, the tool's epoch and time, and a nonce issued by the tool. For $\iota$, the wallet picks the next unused counter $j$ in its own part of $[\ell,k)$ for $(\theta,e)$ and computes $\nym=B_{\theta,e}^{1/(s+j+1)}$ with $B_{\theta,e}=H(\theta,e)$; the base depends only on public tool data and the epoch, never on the tree's secrets $s$ and $t$. It proves, with Fiat--Shamir context $\ctx=\mathrm{SHA\mbox{-}256}(\iota)$ and against the accumulator value $V$ it holds:
\begin{align*}
\mathsf{PoK}\{&\text{signature on } m \text{ with } b_c=1,\\
&\nym^{s+j+1}=B_{\theta,e},\ \ell\le j<k,\ \beta\ge\mathit{amt},\ \tau\ge \mathit{tm},\\
&\mathit{id},\mathit{anc}_{1..4}\notin\mathsf{Acc}_V,\\
&[\,E=(g^{r_1},g^{t}X^{r_1},g^{r_2},g^{\mathit{id}}X^{r_2})\,]\}.
\end{align*}
$\ell\le j<k$ is proved by two ranges, $j-\ell+1\in[1,2^{40})$ and $k-j-1\in[0,2^{40})$. The escrow part $E$ is included when the tool requires it. The output is $(V,\nym,[E],\pi)$; $j$ stays hidden.

\paragraph*{Verify} The tool performs these checks in order:
(1) rebuilds $\iota$ from the raw request, its identity, its clock and the nonce it issued; the nonce is consumed; (2) checks that $e$ is its current epoch and the time is within tolerance; (3) checks that $V$ is the current accumulator value or one of the $W$ previous ones; (4) reserves $(e,\nym)$ in its seen-set; (5) verifies $\pi$, and releases the reservation on failure.
All failures return the same error.

\paragraph*{Delegate} The issuer sends a nonce. The parent chooses the child's attributes, normally with a sub-interval $[\ell',k')$ disjoint from its own part and from its other children. It commits to all child messages except $\mathit{id}'$ and the version ($C_p$), computes a delegation pseudonym $\nym_D=D_e^{1/(s+j+1)}$ for a counter in its own part, and proves, with context $H(\text{nonce}\|C_p)$:
its signature and the non-membership of its five identifiers; by witness equalities: $s'=s$, $t'=t$, $\mathit{anc}'=(\mathit{id},\mathit{anc}_1,\mathit{anc}_2,\mathit{anc}_3)$; $\delta'=\delta+1\le4$; $\beta'\le\beta$, $\tau'\le\tau$, $\ell'\ge\ell$, $k'\le k$, each as a difference in $[1,2^{40})$, together with child values in $[0,2^{40})$; for each class, $b_c-b'_c+1\in[1,2]$; the pseudonym.
The issuer checks that $e$ is its own current epoch, verifies, rejects a seen $(e,\nym_D)$, draws a fresh $\mathit{id}'$ and blind-signs. The parent unblinds and hands the mandate to the child. The issuer learns $\mathit{id}'$, the epoch and the time of the delegation, and nothing else.

The ranges on the child values make attenuation hold over the integers; with differences alone, a child value could wrap around $q$. The scope rule needs no bit proof: a presentation reveals $b_c=1$ exactly, and by induction no descendant of a $0$ (or negative) value can reach $1$.

\paragraph*{Revoke} The log accepts a revocation of $X\neq0$ only with an anonymous proof made with the credential of the live mandate $X$ itself, which only its holder and the parent that delegated it keep (so no descendant revokes an ancestor), or with signatures of two escrow members; identifier $0$, the padding for missing ancestors, is never revocable. The issuer then appends $X$ to the accumulator and publishes the new value with its update data and a spot check (witnesses, under both values, for a point hashed from the update), which anyone can verify; holders refresh their five witnesses from public data, and a wallet that cannot refresh stops presenting. Because each mandate carries its ancestors' identifiers, revoking one identifier stops a subtree. Presentations in flight survive unrelated revocations through a time-bounded window: a tool accepts the previous value only for $\Delta$ seconds after its successor reaches it, so revocation takes effect after $\Delta$ plus propagation, and with batched updates at most two values coexist. A count-based window or a wallet retry are alternatives (E4); the retry links the two attempts at one tool.

\paragraph*{Escrow and Judge} Committee members hold shares of $x_{\mathcal{E}}$ from the DKG of Gennaro et al.~\cite{gennaro2007dkg}. Member $i$ returns $E_1^{x_i}$ and $E_3^{x_i}$ with Chaum--Pedersen proofs~\cite{chaum1992dleq}. Two shares give $g^{t}$ and $g^{\mathit{id}}$. $\mathsf{Judge}$ maps $g^{t}$ to a logged binding, re-verifies the principal's signature on it, and maps $g^{\mathit{id}}$ to the issued mandate.

\section{Security}\label{sec:security}
The supplementary material proves the six properties of Table~\ref{tab:props}. Four points deserve mention.
\begin{itemize}
\item \emph{Arithmetic}: the bound check for $[\ell,h)$ proves $m-\ell$ and $h-1-m$ in $[0,2^{64})$; with $\ell,h<2^{41}<2^{65}<q$, every attenuation, interval and counter statement holds over the integers.
\item \emph{Quota}: within one tool identity $\theta$ (origin and tool name) and epoch, the pseudonym is an injective function of the counter, and the counter lies in the root's interval. Epochs are the tools' and the issuer's, so they cannot be invented. Sub-intervals give children shares of the quota that add up to the root's without coordination. A parent that hands out overlapping intervals only causes rejections inside its own tree.
\item \emph{Issuer}: unlinkability holds against a corrupted issuer that publishes a consistent log, and non-frameability as long as registrations are in that log rather than in the issuer's database. Because the log accepts revocations only from a mandate's holder or parent or from two committee members, the issuer cannot revoke a root to watch which traffic stops; it can still hand one wallet an invalid initial witness, which that wallet detects (tampered update data fail the public spot check).
\item \emph{Trees}: one tree shares $s$ and $t$. A delegate that runs its own wallet can link that tree's calls, and its actions open to the principal. It never receives $u_P$, so it cannot obtain roots or affect the principal's other trees.
\item \emph{Concrete security}: the unlinkability reduction loses $n_h^2Q$ for $n_h$ honest trees and rests on $(Q+1)$-DDHI (about $2^{117.5}$ under Cheon's attack for $Q=2^{20}$), so it gives $2^{80}$ only for a few hundred honest trees per issuer key: deployments should keep $Q$ small and rotate the key; this limits the proof, it is not an attack.
\end{itemize}

\section{Implementation}\label{sec:impl}
{\sloppy The prototype is a Rust crate on \texttt{proof\_system}~0.34, \texttt{bbs\_plus}~0.25, \texttt{vb\_accumulator}~0.29 and \texttt{bulletproofs\_plus\_plus}~0.9~\cite{dockcrypto}, over BLS12-381. A PyO3 module exposes two roles:
an \texttt{Authority}, which plays the issuer, the registry, the escrow and the wallets of the test harness, and allocates disjoint intervals; a \texttt{Verifier}, built from public data only, which issues nonces and checks the epoch and the time.
The MCP binding carries $\{V,\nym,[E],\pi\}$ in the request's \texttt{\_meta} field~\cite{mcpspec}. Each server exposes a \texttt{challenge} tool that returns a single-use nonce and the server's time. The server rebuilds $\iota$ from its origin and tool name, the RFC~8785 canonical JSON of the raw call~\cite{rfc8785}, and its descriptor (class, amount, currency). The \texttt{challenge} tool is called by the client, not offered to the language model. The binding, which we have not yet proposed upstream, is outside the core MCP authorization specification, which requires an OAuth bearer token on every HTTP request and forbids servers to accept other tokens~\cite{mcpauth}; it fits the specification's optional authorization extensions, alongside a per-server access token issued to the domain and shared by all its principals, which identifies the domain, not the user (supplementary material). The same envelope fits A2A message metadata~\cite{a2aspec}, which we did not evaluate. The implementation was written with an AI assistant from the construction above and reviewed by the first author (see Acknowledgment); it has not been audited independently.\par}
